\chapter{Arithmetic in Pascal's triangle}
\chaptertag{Binomial coefficients and divisibility}

\begin{definition}
For $0\le k\le n$,
\[
\binom nk=\frac{n!}{k!(n-k)!}.
\]
It counts the $k$-element subsets of an $n$-element set.
\end{definition}

The counting meaning proves integrality without asking the fraction to look integral. It also makes identities visible.

\begin{proposition}[Pascal and symmetry]
\[
\binom nk=\binom n{n-k},\qquad
\binom nk=\binom{n-1}{k-1}+\binom{n-1}{k}.
\]
\end{proposition}

For Pascal's identity, distinguish $k$-subsets that contain a chosen element from those that do not.

\begin{center}
\begin{tikzpicture}[x=.92cm,y=.70cm]
\foreach \n in {0,...,6}{
  \foreach \k in {0,...,\n}{
    \pgfmathtruncatemacro{\val}{factorial(\n)/(factorial(\k)*factorial(\n-\k))}
    \node[font=\small] (p\n-\k) at ({\k-\n/2},-\n) {$\val$};
  }
}
\draw[teal,thick,rounded corners] ($(p4-1)+(-.3,.28)$) rectangle ($(p4-1)+(.3,-.28)$);
\draw[teal,thick,rounded corners] ($(p4-2)+(-.3,.28)$) rectangle ($(p4-2)+(.3,-.28)$);
\draw[coral,thick,rounded corners] ($(p5-2)+(-.34,.28)$) rectangle ($(p5-2)+(.34,-.28)$);
\draw[-{Stealth},teal] (p4-1)--(p5-2);
\draw[-{Stealth},teal] (p4-2)--(p5-2);
\node[coral,right=15mm of p5-2] {$4+6=10$};
\end{tikzpicture}
\end{center}

\section{Three ways to see the binomial theorem}

\begin{theorem}[Binomial theorem]
\[
(x+y)^n=\sum_{k=0}^n\binom nk x^{n-k}y^k.
\]
\end{theorem}

In expanding $(x+y)^n$, choose which $k$ of the $n$ factors contribute $y$. This proof also explains why the coefficients are integers and why their sum is $2^n$.

\begin{proposition}[Vandermonde]
For nonnegative integers $r,s,n$,
\[
\sum_k\binom rk\binom s{n-k}=\binom{r+s}{n}.
\]
\end{proposition}

Count $n$-person committees from two groups: either all at once, or according to how many come from the first group.

\section{Prime rows}

\begin{theorem}
If $p$ is prime and $1\le k\le p-1$, then
\[
p\divides\binom pk.
\]
\end{theorem}

\begin{proof}
The identity $k\binom pk=p\binom{p-1}{k-1}$ shows $p$ divides $k\binom pk$. Since $p\ndivides k$, Euclid's lemma gives $p\divides\binom pk$.
\end{proof}

Thus the $p$th row of Pascal's triangle becomes $1,0,\dots,0,1$ modulo $p$, which again yields
\[
(x+y)^p\congr x^p+y^p\pmod p.
\]

\section{Measuring divisibility one prime at a time}

\begin{definition}
For nonzero integer $m$, $v_p(m)$ is the exponent of $p$ in its prime factorization. Thus $v_2(40)=3$.
\end{definition}

Valuations turn multiplication into addition:
\[
v_p(ab)=v_p(a)+v_p(b),\qquad
v_p\!\left(\frac ab\right)=v_p(a)-v_p(b)
\]
when the quotient is an integer.

\begin{theorem}[Legendre]
\[
v_p(n!)=\floor{\frac np}+\floor{\frac n{p^2}}+\floor{\frac n{p^3}}+\cdots.
\]
\end{theorem}

\begin{proof}
Among $1,\dots,n$, exactly $\floor{n/p}$ contribute at least one $p$; $\floor{n/p^2}$ contribute one additional $p$; and so on.
\end{proof}

Consequently,
\[
v_p\binom nk=\sum_{j\ge1}\left(\floor{\frac n{p^j}}-\floor{\frac k{p^j}}-\floor{\frac{n-k}{p^j}}\right).
\]
Each parenthesis is $0$ or $1$: it records whether adding $k$ and $n-k$ creates a carry at the $p^j$ place.

\begin{theorem}[Kummer's carry rule]
$v_p\binom nk$ equals the number of carries when $k$ and $n-k$ are added in base $p$.
\end{theorem}

The formula above is already the proof once carries are read from right to left. It gives a remarkable visual test: $\binom nk$ is not divisible by $p$ precisely when subtracting $k$ from $n$ in base $p$ needs no borrowing.

\begin{example}
To find $v_2\binom{10}{3}$, add $3=0011_2$ and $7=0111_2$:
\[
\begin{array}{r}0011\\+0111\\\hline1010\end{array}
\]
There are three carries, so $v_2\binom{10}{3}=3$. Indeed, $\binom{10}{3}=120=2^3\cdot15$.
\end{example}

\begin{labbox}
Color entry $\binom nk$ when it is nonzero modulo $2$, then modulo $3$. Zoom out. The self-similar patterns are the visual shadow of base-$p$ digits and lead to Lucas's theorem.
\end{labbox}

\begin{exercisebox}
\begin{enumerate}
  \item Prove $\sum_k\binom nk=2^n$ and $\sum_k(-1)^k\binom nk=0$ for $n>0$. \exlevel{Fluency}
  \item Compute $v_2(100!)$, $v_5(100!)$, and the number of trailing zeros of $100!$. \exlevel{Fluency}
  \item Give a counting proof of $k\binom nk=n\binom{n-1}{k-1}$. \exlevel{Reasoning}
  \item Use Vandermonde's identity to prove $\sum_k\binom nk^2=\binom{2n}{n}$. \exlevel{Reasoning}
  \item Determine all $n$ for which every interior entry of row $n$ is even. Use Kummer's rule. \exlevel{Reasoning}
  \item For a fixed prime $p$, explore the first $p^2$ rows of Pascal's triangle modulo $p$. State a digit-based conjecture for when $\binom nk\not\congr0\pmod p$. \exlevel{Exploration}
  \item Prove that $\operatorname{lcm}(1,2,\dots,2n)$ is divisible by $\binom{2n}{n}$ by comparing $p$-adic valuations. \exlevel{Bridge}
\end{enumerate}
\end{exercisebox}

